%% arara directives
% arara: xelatex
% arara: bibtex
% arara: xelatex
% arara: xelatex

%\documentclass{article} % One-column default
\documentclass[twocolumn, switch]{article} % Method A for two-column formatting

\usepackage{preprint}
\usepackage{glossaries}
\usepackage{graphicx}

%% Math packages
\usepackage{amsmath, amsthm, amssymb, amsfonts}

%% Algorithm packages
\usepackage{algorithm}
\usepackage{algpseudocode}

%% Bibliography options
\usepackage[numbers,square]{natbib}
\bibliographystyle{unsrtnat}
%\usepackage{natbib}
%\bibliographystyle{Geology}

%% General packages
\usepackage[utf8]{inputenc}	% allow utf-8 input
\usepackage[T1]{fontenc}	% use 8-bit T1 fonts
\usepackage{xcolor}		% colors for hyperlinks
\usepackage[colorlinks = true,
            linkcolor = purple,
            urlcolor  = blue,
            citecolor = cyan,
            anchorcolor = black]{hyperref}	% Color links to references, figures, etc.
\usepackage{booktabs} 		% professional-quality tables
\usepackage{nicefrac}		% compact symbols for 1/2, etc.
\usepackage{microtype}		% microtypography
\usepackage{lineno}		% Line numbers
\usepackage{float}			% Allows for figures within multicol
%\usepackage{multicol}		% Multiple columns (Method B)

\usepackage{lipsum}		%  Filler text

 %% Special figure caption options
\usepackage{newfloat}
\DeclareFloatingEnvironment[name={Supplementary Figure}]{suppfigure}
\usepackage{sidecap}
\sidecaptionvpos{figure}{c}

% Section title spacing  options
\usepackage{titlesec}
\titlespacing\section{0pt}{12pt plus 3pt minus 3pt}{1pt plus 1pt minus 1pt}
\titlespacing\subsection{0pt}{10pt plus 3pt minus 3pt}{1pt plus 1pt minus 1pt}
\titlespacing\subsection{0pt}{8pt plus 3pt minus 3pt}{1pt plus 1pt minus 1pt}

% ORCiD insertion
\usepackage{tikz,xcolor,hyperref}

\definecolor{lime}{HTML}{A6CE39}
\DeclareRobustCommand{\orcidicon}{
	\begin{tikzpicture}
	\draw[lime, fill=lime] (0,0) 
	circle [radius=0.16] 
	node[white] {{\fontfamily{qag}\selectfont \tiny ID}};
	\draw[white, fill=white] (-0.0625,0.095) 
	circle [radius=0.007];
	\end{tikzpicture}
	\hspace{-2mm}
}
\foreach \x in {A, ..., Z}{\expandafter\xdef\csname orcid\x\endcsname{\noexpand\href{https://orcid.org/\csname orcidauthor\x\endcsname}
			{\noexpand\orcidicon}}
}
% Define the ORCID iD command for each author separately. Here done for two authors.
\newcommand{\orcidauthorA}{0009-0004-5890-4819}
\newcommand{\orcidauthorB}{0009-0006-4869-9705}

%%%%%%%%%%%%%%%%   Title   %%%%%%%%%%%%%%%%
\title{A Fast GRASP Metaheuristic for the Trigger Arc TSP}

% Add watermark with submission status
\usepackage{xwatermark}
% Left watermark
\newwatermark[firstpage,color=gray!60,angle=90,scale=0.32, xpos=-4.05in,ypos=0]{\href{https://doi.org/}{\color{gray}{Publication doi}}}
% Right watermark
\newwatermark[firstpage,color=gray!60,angle=90,scale=0.32, xpos=3.9in,ypos=0]{\href{https://doi.org/}{\color{gray}{Preprint doi}}}
% Bottom watermark
\newwatermark[firstpage,color=gray!90,angle=0,scale=0.28, xpos=0in,ypos=-5in]{*correspondence: \texttt{joan.salvasoler@h2o.ai}}

%%%%%%%%%%%%%%%  Author list  %%%%%%%%%%%%%%%
\usepackage{authblk}
\renewcommand*{\Authfont}{\bfseries}
\author[1]{Joan Salva Soler\orcidA{}}
\author[2]{Grégoire de Lambertye\orcidB{}}

\affil[1]{H2O.ai, TU Wien}
\affil[2]{TU Wien}

%%%%%%%%%%%%%%    Front matter    %%%%%%%%%%%%%%
\begin{document}

\twocolumn[ % Method A for two-column formatting
  \begin{@twocolumnfalse} % Method A for two-column formatting
  
\maketitle

\begin{abstract}
 purus elit, vestibulum ut, placerat ac, adipiscing vitae, felis.
 Curabitur dictum gravida mauris. Nam arcu libero, nonummy
 eget, consectetuer id, vulputate a, magna. Donec vehicula au
gue eu neque. Pellentesque habitant morbi tristique senectus
 et netus et malesuada fames ac turpis egestas. Mauris ut leo.
 Cras viverra metus rhoncus sem. Nulla et lectus vestibulum
 urna fringilla ultrices. Phasellus eu tellus sit amet tortor gravida
 placerat. Integer sapien est, iaculis in, pretium quis, viverra ac,
 nunc. Praesent eget sem vel leo ultrices bibendum. Aenean

\end{abstract}

\keywords{TSP \and TA TSP \and GRASP \and Local Search \and Trigger Arc TSP \and Optimisation} % (optional)

\vspace{0.35cm}

  \end{@twocolumnfalse} % Method A for two-column formatting
] % Method A for two-column formatting

%\begin{multicols}{2} % Method B for two-column formatting (doesn't play well with line numbers), comment out if using method A


%%%%%%%%%%%%%%%  Main text   %%%%%%%%%%%%%%%
% \linenumbers
\newacronym{tsp}{TSP}{Traveling Salesman Problem}
\newacronym{tatsp}{TA-TSP}{Trigger Arc TSP}
\newacronym{grasp}{GRASP}{Greedy Randomized Adaptive Search Procedures}
\newacronym{mip}{MIP}{Mixed-Integer Programming}

\section{Introduction}
The \gls{tsp} is a foundational problem in combinatorial optimization. In its traditional form, 
it assumes that the cost of traveling between cities remains constant.
Such assumption fails to capture the dynamic nature of many real-world scenarios. 
The cost of traversing a path can change based on previous actions or on the state of the system. 
One such compelling application arises in warehouses equipped with compactable storage systems. 
These systems optimize space by using mobile shelving units that move on rails. 
Accessing an item in one aisle may require moving shelves, which in turn blocks or alters the cost and time required to access other aisles. 
This path-dependent cost structure requires more sophisticated modeling than what the classical \gls{tsp} offers. 

To address such behaviors, ~\citet{Cerrone} introduced the \gls{tatsp}. 
The \gls{tatsp} is a variant of the \gls{tsp} where the cost of an arc can be dynamically altered by traversing another designated "trigger" arc earlier in the tour. 
Specifically, the problem defines relations between pairs of arcs, where one is a trigger and the other a target. 
When a trigger arc is traversed, it can modify the cost of its associated target arc. 
If multiple trigger arcs for the same target arc are traversed, only the last one encountered before the target arc has an effect. 
This novel formulation provides a powerful tool for modeling problems with state-dependent costs. 

This paper is motivated by the Metaheuristics Summer School (MESS) 2024 competition, which challenged participants to develop effective metaheuristic solutions for the \gls{tatsp}. 
Our work, which secured a top-three position in the competition, proposes a scalable and fast \gls{grasp}-based approach to find high-quality solutions for this complex problem. 
We present a construction heuristic to generate diverse and feasible initial solutions, followed by a local search procedure to iteratively improve them. 
This paper details our methodology, experimental setup, and the results of applying our \gls{grasp} algorithm to a set of benchmark instances, demonstrating its effectiveness and efficiency in tackling the \gls{tatsp}. 

\section{Related Work}
\label{sec:related_work}

\subsection{Variations on the \gls{tsp}}

The \gls{tsp} is fundamentally about discovering the most efficient route that connects a given set of nodes~\cite{Reinelt2003}. 
Its NP-hard nature makes exact solutions for large instances computationally infeasible, necessitating the use of heuristics and metaheuristics.
Their utility becomes even more pronounced when dealing with the intricacies of dynamic TSP variants. These variants introduce 
real-world complexities, including problems on temporal graphs where edge access is time-sensitive, TSPs where visiting a node 
has the immediate consequence of deleting specific edges, and scenarios where the cost structure of a tour is directly affected 
by the historical sequence of visited locations. A prime example is the the Node Weight Dependent TSP (W-TSP), where the cost of a 
tour increases based on the weights of previously visited nodes~\cite{Bossek2020}.

The \gls{tatsp}, introduced by~\citet{Cerrone}, is a variant that introduces dynamic arc costs influenced by the solution's path. 
In the \gls{tatsp}, the cost of traversing a specific arc can be reset if a predefined "trigger arc" is traversed beforehand, 
with only the most recently encountered trigger arc affecting the target arc's cost. In their introduction paper, ~\citet{Cerrone} 
propose an Integer Linear Programming (ILP) model for the \gls{tatsp}. Their approach provides a precise mathematical formulation 
but is limited in scalability. This inherent complexity and the practical need for fast solutions for large-scale instances underscore
 the need for robust metaheuristic approaches.

\subsection{Metaheuristics for the TSP}
\label{sec:metaheuristics_tsp}

Metaheuristics are a class of optimization algorithms that provide approximate solutions to complex problems by exploring and exploiting the search space efficiently~\cite{Gendreau2010}. Due to the NP-hard nature of the \gls{tsp}, metaheuristics have been extensively explored and applied to find high-quality solutions within reasonable computational time~\cite{toaza2023}. 
The computational difficulty of the \gls{tsp} has led to the development of a wide array of metaheuristic approaches. Prominent examples include Genetic Algorithms (GAs), which are inspired by natural evolution; Tabu Search (TS), which uses memory structures to escape local optima; Simulated Annealing (SA), a probabilistic method that accepts worse solutions with a decreasing probability; and Ant Colony Optimization (ACO), which mimics the foraging behavior of ants. These methods have been extensively studied and adapted for the \gls{tsp} and its variants, demonstrating the critical role of metaheuristics in addressing NP-hard optimization problems. The success of these diverse approaches provides a strong precedent for applying a metaheuristic framework to the even more complex \gls{tatsp}.

\subsection{The GRASP approach}

This work is based on the \gls{grasp} metaheuristic, introduced by Feo and Resende~\cite{Feo1995}, which is a two-phase approach that combines a construction heuristic and a local search procedure.
The construction phase leverages a greedy randomized strategy to build a diverse set of feasible solutions. Following this, the local search phase systematically explores the neighborhood of each constructed solution to find a local optimum with respect to a set of neighborhood structures.
The best solution found across all iterations is returned.The two-phase structure of \gls{grasp} inherently balances diversification (exploration) and intensification (exploitation) of the search space, which is advantageous for complex fitness landscapes.
Due to its robust performance and adaptability, \gls{grasp} has been successfully applied to the classical \gls{tsp}, often achieving results competitive with highly specialized heuristics~\cite{Oliveira2004}. 
Its effectiveness can be further enhanced through advanced techniques. Reactive GRASP, for instance, dynamically adjusts its parameters during the search to adapt to the problem's characteristics~\cite{Prais2000}. 
Path-Relinking is an intensification strategy that explores trajectories between elite solutions to find improved ones~\cite{Resende2019}. The consistent success of GRASP and its variants in related routing problems makes it a compelling choice for addressing the \gls{tatsp}.

\section{Methodology}

\subsection{Problem Definition}

The \gls{tatsp} is formally defined on a directed graph $G = (N, A)$, where $N$ is the set of nodes and $A$ is the set of directed arcs. Node $0 \in N$ is designated as the starting node (depot). The problem incorporates the following key elements:

\begin{itemize}
\item \textbf{Base arc costs}: Let $c(a) \in \mathbb{R}^+$ for all $a \in A$ represent the base cost of traversing arc $a$.
\item \textbf{Relations}: For each arc $a_j \in A$, a set of relations $R_a = \{(a_i, a_j) \mid a_i \in A\}$ is defined, where each relation connects a trigger arc $a_i$ to the target arc $a_j$.
\item \textbf{Relation costs}: Let $c(r) \in \mathbb{R}^+$ for all $r \in R$ represent the cost associated with relation $r$ when active, where $R = \bigcup_{a \in A} R_a$ is the set of all relations.
\end{itemize}

Given a Hamiltonian cycle $T = (a_i, a_j, \ldots, a_{|N|})$ starting at node $0$, a relation $r = (a_i, a_j) \in R$ is \textbf{active} if and only if:
\begin{enumerate}
\item Both arcs $a_i$ and $a_j$ belong to the tour $T$
\item Arc $a_i$ precedes arc $a_j$ in the traversal sequence
\item There is no other relation $r' = (a_k, a_j) \in R_{a_j}$ such that $a_k$ follows $a_i$ in $T$
\end{enumerate}

The cost of traversing an arc $a \in A$ in tour $T$ is determined as:
\begin{equation}
\text{cost}(a, T) = \begin{cases}
c(r) & \text{if relation } r \in R_a \text{ is active} \\
c(a) & \text{if no relation in } R_a \text{ is active}
\end{cases}
\end{equation}

The objective of the \gls{tatsp} is to find a Hamiltonian cycle $T^*$ that minimizes the total tour cost:
\begin{equation}
T^* = \arg\min_{T} \sum_{a \in T} \text{c}(a)
\end{equation}

One observation is that the \gls{tatsp} generalizes the \gls{tsp}. If the relation set within the TA-TSP is empty, 
the problem formulation simplifies directly to that of the traditional TSP.
The two problems share the same set of feasible solutions, which is the set of all Hamiltonian cycles.
However, there is no natural reduction from the \gls{tatsp} to the \gls{tsp} because of the path-dependency of the cost function.

For modeling purposes, it is useful to reformulate the cost function. Instead of viewing the cost of a relation as a replacement for the base arc cost, we can consider it as an additional cost (or discount). For each relation $r = (a_i, a)$ with an associated cost $c(r)$, we can define a delta cost, $\Delta(r) = c(r) - c(a)$. The cost of traversing an arc $a$ in a tour $T$ can then be written as $c(a) + \Delta(r)$ if relation $r$ is active, and $c(a)$ otherwise. This separates the base tour cost from the additional costs incurred by active relations, which is a convenient perspective for developing heuristics and mathematical models.

\subsection{The \gls{grasp} approach}

For solving the \gls{tatsp}, we leverage a classic \gls{grasp} approach as introduces by ~\cite{Feo1995}. 
This framework described in Algorithm~\ref{alg:grasp_general} consists in iterating over these two phases, 
and the best solution found across all iterations is returned. The first phase is the Construction Phase, which builds a feasible
solution using a randomized construction heuristic. The second phase, Local Search Phase, applies improvement procedures to the 
constructed solution until a local optimum is reached. 

\begin{algorithm}
\caption{General GRASP Framework}
\label{alg:grasp_general}
\begin{algorithmic}[1]
\State $S^* \leftarrow \emptyset$ \Comment{Best solution found}
\State $f^* \leftarrow \infty$ \Comment{Best objective value}
\For{$i = 1$ to $\text{MaxIterations}$}
    \State $S \leftarrow \text{RandomizedConstruction}()$
    \State $S \leftarrow \text{LocalSearch}(S)$
    \If{$f(S) < f^*$}
        \State $S^* \leftarrow S$
        \State $f^* \leftarrow f(S)$
    \EndIf
\EndFor
\State \Return $S^*$
\end{algorithmic}
\end{algorithm}

As a local search procedure, we use a multi-neighborhood local search procedure with a first-improvement policy.
The algorithm is presented in Algorithm~\ref{alg:local_search}.

\begin{algorithm}
\caption{Multi-Neighborhood Local Search}
\label{alg:local_search}
\begin{algorithmic}[1]
\State \textbf{Input:} Tour $T$, Neighborhoods $\mathcal{N} = \{\mathcal{N}_1, \mathcal{N}_2, \mathcal{N}_3\}$
\State \textbf{Output:} Locally optimal tour $T'$
\Repeat
    \State $\text{improved} \leftarrow \text{false}$
    \For{each neighborhood $\mathcal{N}_i \in \mathcal{N}$}
        \For{each move $m \in \mathcal{N}_i(T)$}
            \State $T' \leftarrow \text{ApplyMove}(T, m)$
            \State $\Delta \leftarrow \text{EvaluateMoveCost}(T, m)$ \Comment{Compute cost change with trigger effects}
            \If{$\Delta < 0$}
                \State $T \leftarrow T'$
                \State $\text{improved} \leftarrow \text{true}$
                \State \textbf{break} \Comment{First improvement strategy}
            \EndIf
        \EndFor
        \If{$\text{improved}$}
            \State \textbf{break} \Comment{Restart with first neighborhood}
        \EndIf
    \EndFor
\Until{$\text{improved} = \text{false}$}
\State \Return $T$
\end{algorithmic}
\end{algorithm}

\subsection{Construction Heuristics}
This section defines the construction heuristics benchmarked in this study.

\paragraph{Simple Randomized Construction}
This heuristic serves as a baseline method for generating feasible solutions. It constructs a tour iteratively, starting from the depot (node 0). At each step, it identifies all unvisited neighbors of the last node added to the tour. From this set of feasible next nodes, one is selected uniformly at random and appended to the tour. This process continues until all nodes have been visited. The final node is only considered feasible if it has an arc connecting back to the depot, thus ensuring a valid Hamiltonian cycle. This method does not consider any cost information during construction, relying purely on random selection to explore the solution space.

\paragraph{Randomized Greedy Construction}
This method is a semi-greedy heuristic, similar to the simple randomized approach, it builds a solution one node at a time. However, the choice of the next node is guided by cost information. At each step, for every feasible successor node, the exact incremental cost of adding the corresponding arc to the partial tour is calculated. This evaluation is non-trivial, as it must account for any changes in the total tour cost due to the activation of trigger-target arc relations.

A Restricted Candidate List (RCL) is then formed, containing the most promising arcs based on their incremental cost. The size of the RCL is controlled by a randomization parameter $\alpha \in [0, 1]$. $\alpha$ correesponds to the share of neighbors added to the short list. An arc is then selected uniformly at random from the RCL and the corresponding node is added to the tour. A value of $\alpha = 0$ results in a purely greedy construction, where the best arc is always chosen, while $\alpha = 1$ allows for a random selection from all feasible arcs.
  
\paragraph{\textit{MIP-based Random Perturbation}}
This heuristic leverages a full \gls{tsp} solver by transforming the \gls{tatsp} into a classic \gls{tsp}, thereby ignoring all trigger-target relations.
The core idea is to perturb the original edge costs $c_{ij}$ to generate a new \gls{tsp} instance.
This instance is then solved to optimality or near-optimality, and its resulting tour (or a pool of high-quality tours) is evaluated using the true \gls{tatsp} objective function.
In the \textit{Additive Perturbation} approach, the costs are modified using an additive noise factor controlled by a parameter $\alpha$. The new cost $c'_{ij}$ is calculated as $c'_{ij} = c_{ij} + \alpha \cdot U(-1, 1)$, where $U(-1, 1)$ is a random variable from a uniform distribution between -1 and 1. Alternatively, the \textit{Multiplicative Perturbation} method scales the costs by a random factor controlled by a parameter $\beta$, resulting in $c'_{ij} = c_{ij} \cdot (\beta \cdot U(0, 1))$.
  
\paragraph{\textit{MIP-based Biased Perturbation}}
This method attempts to incorporate relational information into the \gls{tsp} model by perturbing arc costs based on an estimated probability of relation activation. The process begins by generating a random permutation of nodes, $P$, which serves as a prior for estimating arc usage. The probability of an arc $(i, j)$ being used is assumed to be inversely proportional to the cyclic distance between nodes $i$ and $j$ in the prior permutation $P$. This is defined as $p_{ij} = 1/d_{ij}$, where $d_{ij} = \min(|pos(i)-pos(j)|, |N| - |pos(i)-pos(j)|)$. For a relation $r$ from a trigger arc $a=(a_i, a_j)$ to a target arc $b=(b_1, b_2)$, the activation probability $p_r$ is estimated by considering the usage probabilities of both arcs and their proximity. It is defined as $p_r = p_{a_ia_j} \cdot p_{b_1b_2} / d_{a_jb_1}^\beta$, where $d_{t_2g_1}$ is the cyclic distance between the trigger's endpoint and the target's startpoint, and $\beta$ is a tunable parameter. The costs of the arcs in the base \gls{tsp} instance are modified to reflect the expected impact of the relations. For each relation $r$ between trigger arc $b$ and target arc $a$, a penalty proportional to the relation's activation probability is added to the costs of both arcs. The new cost of an arc $k \in \{a,b\}$ is updated as $c'(k) = c(k) + \alpha \cdot p_r \cdot c_r$, where $c_r$ is the cost when relation $r$ is active, and $\alpha$ is a parameter controlling the penalty's magnitude. Finally, the modified \gls{tsp} instance is solved, and the resulting tours are evaluated on the original \gls{tatsp} instance to find the best solution.

\subsection{Local search}
The local search phase is designed to iteratively improve a solution by exploring its neighborhood. We employ a multi-neighborhood variable search procedure that systematically applies a set of local search operators. The search follows a first-improvement strategy: as soon as a move that reduces the total tour cost is found, the move is applied, and the search is restarted from the beginning of the neighborhood list. This process continues until no further improvements can be found in any of the defined neighborhoods, at which point the solution is considered locally optimal. The neighborhood structures used are 2-Opt, Swap, and Relocate.

\subsubsection{Neighborhood Structures}
The neighborhoods we employ are standard in the \gls{tsp} literature:

\paragraph{2-Opt} The 2-Opt move is a powerful neighborhood operator for path-improvement~\cite{Croes1958}. Given a tour, it works by removing two non-adjacent arcs, $(v_i, v_{i+1})$ and $(v_j, v_{j+1})$, and reconnecting the two resulting paths in the only other possible way. This is equivalent to reconnecting them with arcs $(v_i, v_j)$ and $(v_{i+1}, v_{j+1})$, which has the effect of reversing the segment of the tour between nodes $v_{i+1}$ and $v_j$.

\paragraph{Swap} The Swap (or Exchange) move consists of exchanging the positions of two nodes, $v_i$ and $v_j$, in the tour. The depot node is not considered for this operation. This move alters four arcs in the tour: the two arcs entering and the two arcs leaving the swapped nodes.

\paragraph{Relocate} The Relocate (or Insertion) move selects a node $v_i$ and removes it from its current position in the tour. It is then re-inserted into a different position, between two adjacent nodes $(v_j, v_{j+1})$.

\subsection{Delta Evaluation}
A critical aspect of any iterative improvement algorithm is the evaluation of moves. The efficiency of construction heuristics and local search often relies on fast "delta evaluation"—calculating the change in objective value without recomputing the entire solution's cost.

For our sequential construction heuristics, such as the Randomized Greedy method, we successfully implement a form of delta evaluation. When considering adding a new node to a partial tour, we only need to compute the cost of the new arc and any trigger effects it may be subject to from arcs already in the partial tour. This is possible because the partial tour is fixed and extended in a single direction, so the trigger states of existing arcs are not altered.

However, this approach is not feasible for the local search phase. Standard neighborhood moves, like 2-Opt or Swap, can drastically alter the sequence of the tour. A single swap can change which trigger is the "last" one to activate a target arc located much later in the tour. This creates a cascade effect, where a local change can lead to non-local cost modifications across the entire tour. A true delta evaluation would require a complex procedure to track and re-evaluate the effective trigger for every potentially affected arc, defeating the performance benefits of the technique. For correctness and simplicity, our local search implementation re-evaluates the entire tour cost after every move.

\subsection{MIP Model}
For benchmarking purposes, we developed a MIP model for the \gls{tatsp}, which shares similarities with the ILP model proposed by~\citet{Cerrone}. We use the commercial solver Gurobi \cite{gurobi}to solve the MIP model.
The full model is presented in the Appendix.


\section{Computational Experiments}
This section presents the computational experiments conducted to evaluate the performance of the proposed \gls{grasp} algorithm.

\subsection{Datasets}


\begin{table*}[!htpb]
    \centering
    \caption{Dataset Characteristics}
    \label{tab:dataset_characteristics}
    \resizebox{\textwidth}{!}{%
        \begin{tabular}{|l|l|l|l|l|l|l|}
        \hline
        \textbf{Datasets} & \textbf{\# Inst.} & \textbf{\# Nodes (Min/Avg/Max)} & \textbf{\# Arcs (Min/Avg/Max)} & \textbf{\# Rels (Min/Avg/Max)} & \textbf{Avg Arc Cost} & \textbf{Avg Rel Cost} \\
        \hline
        instances\_generic & 180 & 10 / 18 / 25 & 90 / 320 / 600 & 100 / 2092 / 10000 & 2587.19 & 3042.15 \\
        \hline
        instances\_release\_2 & 34 & 18 / 68 / 142 & 90 / 580 / 1562 & 1144 / 47834 / 208020 & 1.00 & 1.08 \\
        \hline
        instances\_release\_1 & 21 & 20 / 40 / 60 & 200 / 1267 / 2700 & 1732 / 556903 / 4527944 & 6.29 & 6.30 \\
        \hline
        \end{tabular}%
    }%
\end{table*}

Three datasets were used to evaluate our \gls{grasp} algorithm and assess its performance across problem having different characteristics.

The first dataset contains 180 synthetic instances generated as described in~\citet{Cerrone}.
Each instance is generated using a random generator that creates complete directed graphs $G(N,A)$ 
with nodes placed within a 5 km square Euclidean space. 
For each pair of nodes $i, j \in N$, directed arcs $(i, j)$ and $(j, i)$ are created and associate with a base cost $c(a)$ 
corresponding to the Euclidean distance between the nodes. The synthetic instances are divided into three scenarios 
(Balanced, Increase, Decrease), with 60 instances per scenario. Each scenario includes instances with $n \in \{10, 15, 20, 25\}$ 
nodes and $|\frac{n}{2}, 2n, 4n, 8n, 16n|$ relations per graph. For each combination of node count and relation count, three 
different instances are created, and all computational results report the average value obtained from these three instances.
Relations $r = (a_i, a_j) \in R$ are generated by randomly selecting pairs of arcs $a_i, a_j \in A$. The cost of each relation $c(r)$ 
is determined by generating a random value in the range $[\frac{c(a_j)}{2}, 2c(a_j)]$.


The three scenarios are characterized as follows:
\begin{itemize}
    \item \textbf{Balanced}: Relations have random costs that can be either greater or less than the base arc cost, in the range $[\frac{c(a_j)}{2}, 2c(a_j)]$
    \item \textbf{Increase}: Relations have random costs greater than the base arc cost, in the range $[c(a_j), 2c(a_j)]$
    \item \textbf{Decrease}: Relations have random costs less than the base arc cost, in the range $[\frac{c(a_j)}{2}, c(a_j)]$
\end{itemize}

Two additional datasets, provided by the competition organizers, offer realistic scenarios with varying problem sizes and complexities. The first of these contains 21 instances, with an average of 580 arcs and 47,834 relations. The second, larger dataset has 34 instances, averaging 1,267 arcs and 556,903 relations. Table \ref{tab:dataset_characteristics} summarizes the key characteristics of all three datasets.

\subsection{Experimental Setup}

Our \gls{grasp} algorithm is implemented in C++20, leveraging modern language features for efficient memory management and performance. The implementation uses the Boost libraries for data structures and the Gurobi optimization suite for solving the MIP-based construction heuristics.

All computational experiments were conducted on a machine equipped with an Apple M1 Pro processor running macOS 24.5.0 (Darwin Kernel Version 24.5.0). The system architecture is ARM64, providing a modern computing environment for evaluating the algorithm's performance.

The complete implementation, along with all instances used in our experiments, is publicly available at \url{https://github.com/jsalvasoler/trigger_arc_tsp}.

\subsection{Metrics}

To evaluate the performance of our \gls{grasp} algorithm, we report the optimality gap, which measures how close our solutions are to the best possible solutions. The gap is calculated as:

\begin{equation}
\text{Gap} = \frac{\text{Solution Cost} - \text{Best Known Cost}}{\text{Best Known Cost}} \times 100\%
\end{equation}

For the competition instances (the first two datasets), we use the best known solutions provided by the competition organizers as the reference point. For the synthetic instances (the third dataset), we use the best dual bound obtained after running the Gurobi MIP solver for one minute as the reference point, since optimal solutions are not known for these instances.

\subsection{Results and Discussion}
\lipsum[11]

\section{Conclusion}
\lipsum[12]

\appendix
\section{MIP Model for the TA-TSP}

This appendix presents the detailed \gls{mip} model for the \gls{tatsp}.

The model uses the following parameters, sets, and decision variables:
\begin{itemize}
    \item $V$: Set of nodes $\{0, 1, \dots, N-1\}$.
    \item $E$: Set of directed arcs $(i,j)$.
    \item $c_{ij}$: Base cost of traversing arc $(i,j)$.
    \item $R_a$: Set of trigger arcs for a target arc $a \in E$.
    \item $r_{ba}$: Additional cost incurred if trigger arc $b$ activates target arc $a$.
    \item $x_{ij} \in \{0, 1\}$: 1 if arc $(i,j)$ is in the tour, 0 otherwise.
    \item $u_i \in [0, N-1]$: Position of node $i$ in the tour. For an arc $a=(i,j)$, $u_a$ denotes $u_i$.
    \item $y_{ba} \in \{0, 1\}$: 1 if relation $(b,a)$ is active, 0 otherwise.
    \item $z_{a_ia_j} \in \{0, 1\}$: 1 if arc $a_i$ precedes arc $a_j$ in the tour, 0 otherwise.
\end{itemize}

\footnotesize

The objective function minimizes the total tour cost:
\begin{equation} \label{eq:obj}
\min \sum_{(i,j) \in E} c_{ij} x_{ij} + \sum_{a \in E} \sum_{(b,a) \in R_a} r_{ba} y_{ba}
\end{equation}

Subject to the following constraints:
\begin{subequations}
\begin{align}
    \sum_{j \in V} x_{ij} &= 1 \quad \forall i \in V \label{eq:flow_out} \\
    \sum_{j \in V} x_{ji} &= 1 \quad \forall i \in V \label{eq:flow_in} \\
    u_i - u_j + N x_{ij} &\leq N-1 \quad \forall (i,j) \in E, j \neq 0 \label{eq:mtz} \\
    u_0 &= 0 \label{eq:start_node} \\
    \sum_{(b,a) \in R_a} y_{ba} &\leq x_a \quad \forall a \in E \label{eq:rel_active_1} \\
    y_{ba} &\leq x_b \quad \forall a \in E, (b,a) \in R_a \label{eq:rel_active_2} \\
    u_b + 1 &\leq u_a + N(1-y_{ba}) \quad \forall a \in E, (b,a) \in R_a \label{eq:rel_precedence} \\
    1 - z_{ab} &\leq \sum_{(c,a) \in R_a} y_{ca} + (1 - x_a) + (1 - x_b) \notag \\
    &\quad \forall a \in E, (b,a) \in R_a \label{eq:rel_active_3} \\
    u_{a_i} &\leq u_{a_j} + (N-1)(1-z_{a_ia_j}) \quad \forall a_i, a_j \in E \label{eq:z_precedence} \\
    y_{ba} &\leq y_{ca} + z_{cb} + z_{ac} + (1-x_c) + (1-x_b) + (1-x_a) \notag \\
    &\quad \forall a \in E, (b,a),(c,a) \in R_a, b \neq c \label{eq:rel_last_trigger}
\end{align}
\end{subequations}

\normalsize

The objective function \eqref{eq:obj} minimizes the total tour cost, composed of the base arc costs and the additional costs from active relations.

Constraints \eqref{eq:flow_out} and \eqref{eq:flow_in} are flow conservation constraints ensuring that each node is visited exactly once.
Constraints \eqref{eq:mtz} and \eqref{eq:start_node} are the Miller-Tucker-Zemlin (MTZ) formulations for subtour elimination, which define the sequence of nodes in the tour, with node 0 as the starting point.

The relation activation logic is modeled by constraints \eqref{eq:rel_active_1} through \eqref{eq:rel_last_trigger}.
Constraint \eqref{eq:rel_active_1} ensures that a relation for a target arc $a$ can only be active if $a$ is part of the tour, and at most one such relation can be active.
Constraint \eqref{eq:rel_active_2} strengthens this by requiring the trigger arc $b$ to also be in the tour for the relation $(b,a)$ to be active.
Constraint \eqref{eq:rel_precedence} states that for a relation $(b,a)$ to be active, the trigger arc $b$ must be traversed before the target arc $a$.
The precedence variable $z_{a_ia_j}$ is defined in constraint \eqref{eq:z_precedence}, where $z_{a_ia_j}=1$ if arc $a_i$ precedes arc $a_j$.
Constraint \eqref{eq:rel_active_3} ensures that if an arc $a$ and one of its potential triggers $b$ are in the tour, and $b$ precedes $a$, then at least one relation for $a$ must be activated.
Finally, constraint \eqref{eq:rel_last_trigger} is the core constraint that models the "last trigger" rule. It ensures that if two triggers $b$ and $c$ for the same target $a$ are in the tour, and $c$ is traversed after $b$ but before $a$, then the relation $(b,a)$ cannot be active.

%%%%%%%%%%%% Supplementary Methods %%%%%%%%%%%%
%\footnotesize
%\section*{Methods}

%%%%%%%%%%%%% Acknowledgements %%%%%%%%%%%%%
\footnotesize
\section*{Acknowledgements}
% thanks to the competition organizers and everything

%%%%%%%%%%%%%%   Bibliography   %%%%%%%%%%%%%%
\normalsize
\bibliography{references}

%%%%%%%%%%%%  Supplementary Figures  %%%%%%%%%%%%
%\clearpage

%%%%%%%%%%%%%%%%   End   %%%%%%%%%%%%%%%%
%\end{multicols}  % Method B for two-column formatting (doesn't play well with line numbers), comment out if using method A
\end{document}
